В общем случае инвариантная мера \textbf{не единственна}, но в следующей теореме, называемой \textit{эргодической}, выделен важный случай, когда она единственна, причем при любом начальном распределении распределения $\pi_n$ сходятся к инвариантному. В данном конечном случае все разумные виды сходимости
мер совпадают, так что можно говорить о сходимости по вариации (по
метрике $\|\mu - \nu\| =
\P|\mu(x_i) - \nu(x_i)|)$ или на каждом элементе из $X$.

\begin{theorem}[Условие эргодичнности конечной цепи Маркова]
    Пусть $X$ конечно и при некотором $n_1$ все элементы матрицы $\P^{n_1}$ положительны. Тогда инвариантная вероятностная мера $\mu$ цепи единственна, причем к ней сходятся меры $\pi_n = \pi_0\P^n$ для всякого начального распределения цепи и $\mu(x_i) > 0$ для всех $i$.
\end{theorem}

\begin{proof}
    Зафиксируем какую-либо инвариантную вероятностную меру $\mu$. Для упрощения обозначений будем считать, что $n_1 = 1$, общий случай аналогичен. Пусть $M = (m_{i,j})$ — матрица, в которой все строки совпадают с $\mu$, т.е. с $(\mu(s_1), \ldots, \mu(s_m))$. Для этой матрицы $\pi M = \mu$ для всех вероятностных мер $\pi$. Найдется такое $\delta \in (0, 1)$, что $p_{ij} \geq \delta m_{ij}$ при всех $i, j$. Положим $\lambda = 1 - \delta$. Тогда $P$ можно представить в виде выпуклой комбинации
    
    \[ P = (1 - \lambda)M + \lambda Q \]
    
    матрицы $M$ и некоторой матрицы $Q = (q_{ij} )$ с $q_{ij} > 0$ и $\sum_j q_{ij} = 1$.
    При таком выборе оказывается, что
    
    \[ P = (1 - \lambda^n)M + \lambda^n Q^n \]
    
    Действительно, непосредственно проверяется, что $M^2 = M$ и $MP = PM = M$.
    Поэтому $QM = M$. Значит, рассуждая по индукции, имеем
    
    \[ P^{n+1} = ((1- \lambda^n) M + \lambda^n Q^n)P=(1-\lambda^n)M+\lambda^n Q^n ((1 - \lambda)M + \lambda Q)=(1-\lambda^{n+1})M+\lambda^{n+1}Q^{n+1}, \]
    
    что завершает индуктивный шаг. Таким образом, с учетом равенства $\pi_0 M - \mu M = 0$ получаем
    
    \[ \pi_0 P^n-\mu = \pi_0 P^n-\mu P^n = (\pi0-\mu)P^n = (\pi0-\mu)((1-\lambda^n)M +\lambda^n Q^n) = \lambda^n (\pi_0 - \mu)Q^n, \]
    
    что стремится к нулю при $n \to \infty$, так как $k(\pi_0 - \mu)Q_{nk} \leq 2$. Если вероятностная мера $\nu$ также инвариантна, то $\nu = P^n \nu \to \mu$, откуда $\nu = \mu$. Наконец, $P^{n_1} \mu(x_i) > 0$ для всех $i$.
\end{proof}

\begin{note}
    Отметим, что условие положительности некоторой степени $P$ и необходимо для справедливости заключения теоремы. В самом деле, если всякая матрица $P^n$ имеет нулевой элемент, то для некоторых $i_0, j_0$ для бесконечной последовательности $n_k$ матрица $P^{n_k}$ имеет нулевой элемент с индексами $i_0, j_0$. Тогда для меры $\pi$, равной нулю на всех $x_i$ c $i \neq i_0$ и $\pi(x_{i_0}) = 1$, получаем $\pi P^{n_k} (x_{j_0}) = 0$, поэтому предел не может быть положительным на всех $x_j$.
\end{note}